\documentclass{article}
\usepackage[hagavi]{azzam}

\begin{document}
\section*{Soal}
\begin{enumerate}
\item Tentukan banyaknya bilangan bulat positif $n \le 1000$ yang merupakan kelipatan 3 atau 4 tetapi bukan kelipatan 5.

\item Tentukan banyaknya pasangan bilangan bulat $(x,y)$ yang memenuhi persamaan $xy = x + y + 2026$.

\item Tentukan banyaknya susunan huruf berbeda yang dapat dibentuk dari kata "OLIMPIADE" sedemikian rupa sehingga semua huruf vokal selalu berdampingan.

\item Tentukan banyaknya titik potong maksimal yang dapat dibentuk oleh 10 lingkaran dan 5 garis lurus di sebuah bidang datar.

\item Tentukan banyaknya bilangan bulat $x$ dengan $1 \le x \le 100$ sedemikian rupa sehingga $x^2 + x + 1$ habis dibagi 3.

\item Tentukan banyaknya solusi bilangan bulat non-negatif dari persamaan $a + b + c + d = 15$.

\item Tentukan banyaknya angka nol berurutan di bagian akhir dari nilai $100!$.

\item Tentukan banyaknya akar real yang memenuhi persamaan $x^4 - 5x^2 + 4 = 0$.

\item Tentukan banyaknya segitiga yang dapat dibentuk menggunakan titik-titik sudut dari sebuah segi-12 beraturan (dodekagon) sedemikian rupa sehingga tidak ada sisi segitiga yang merupakan sisi segi-12 tersebut.

\item Tentukan banyaknya cara 5 siswa dapat mengembalikan 5 buku perpustakaan yang berbeda sedemikian rupa sehingga tidak ada satupun siswa yang menerima buku yang ia pinjam sebelumnya.

\item Tentukan banyaknya faktor pembagi positif dari $10^5$ yang bukan merupakan kelipatan 10.

\item Tentukan banyaknya himpunan bagian dari $\{1, 2, 3, 4, 5, 6, 7, 8\}$ yang mengandung setidaknya satu bilangan genap.

\item Tentukan banyaknya tripel bilangan bulat $(x, y, z)$ yang memenuhi sistem persamaan $x + y = 2$ dan $xy - z^2 = 1$.

\item Tentukan banyaknya diagonal ruang yang terdapat pada sebuah bangun ruang prisma segi-8 beraturan.

\item Tentukan banyaknya rute terpendek di bidang Kartesius dari titik $(0,0)$ menuju $(5,5)$ jika pergerakan hanya diizinkan ke kanan atau ke atas, tanpa melewati titik $(2,2)$.
\end{enumerate}

\section*{Pembahasan}
\begin{enumerate}

\item Misalkan $A$, $B$, dan $C$ berturut-turut adalah himpunan bilangan kelipatan 3, 4, dan 5 yang kurang dari atau sama dengan 1000. Kita mencari $|(A \cup B) \setminus C| = |A \cup B| - |(A \cup B) \cap C|$. 
Banyaknya bilangan kelipatan 3 atau 4 adalah $|A \cup B| = 333 + 250 - 83 = 500$.
Banyaknya bilangan kelipatan 3 atau 4 yang juga merupakan kelipatan 5 adalah kelipatan 15 atau 20, yaitu $|(A \cap C) \cup (B \cap C)| = 66 + 50 - 16 = 100$. 
Hasil akhirnya adalah $500 - 100 = 400$.

\item Persamaan dapat dimanipulasi menjadi $xy - x - y + 1 = 2027$, yang dapat difaktorkan menjadi $(x-1)(y-1) = 2027$. Karena $2027$ adalah bilangan prima, faktor pembaginya di himpunan bilangan bulat hanyalah $1, -1, 2027$, dan $-2027$. Ini memberikan 4 kemungkinan pasangan untuk $(x-1)$ dan $(y-1)$. Jadi, terdapat 4 pasangan $(x,y)$ yang memenuhi.

\item Kata "OLIMPIADE" memiliki 5 huruf vokal (O, I, I, A, E) dan 4 konsonan (L, M, P, D). Jika 5 huruf vokal digabungkan menjadi 1 blok kesatuan, terdapat $1 + 4 = 5$ entitas. Banyaknya cara menyusun 5 entitas ini adalah $5! = 120$. Di dalam blok vokal itu sendiri, banyaknya permutasi dengan memperhitungkan 2 huruf I yang identik adalah $\frac{5!}{2!} = 60$. Total susunan berbeda adalah $120 \times 60 = 7200$.

\item Titik potong maksimal diperoleh jika semua objek saling berpotongan secara optimal. 
Titik potong antar 10 lingkaran: $\binom{10}{2} \times 2 = 90$. 
Titik potong antar 5 garis: $\binom{5}{2} \times 1 = 10$. 
Titik potong antara garis dan lingkaran: $10 \times 5 \times 2 = 100$. 
Total titik potong maksimal adalah $90 + 10 + 100 = 200$.

\item Tinjau persamaan $x^2 + x + 1 \pmod 3$. 
Jika $x \equiv 0 \pmod 3$, maka hasilnya $\equiv 1 \pmod 3$. 
Jika $x \equiv 1 \pmod 3$, maka hasilnya $\equiv 1+1+1 = 3 \equiv 0 \pmod 3$. 
Jika $x \equiv 2 \pmod 3$, maka hasilnya $\equiv 4+2+1 = 7 \equiv 1 \pmod 3$.
Syarat terpenuhi hanya jika $x \equiv 1 \pmod 3$. Dalam rentang $1 \le x \le 100$, nilai yang memenuhi berbentuk deret aritmetika $1, 4, 7, \dots, 100$. Banyaknya adalah $\lfloor \frac{100-1}{3} \rfloor + 1 = 34$.

\item Menggunakan metode \textit{Stars and Bars}, banyaknya solusi bilangan bulat non-negatif dari persamaan $x_1 + x_2 + \dots + x_k = n$ dirumuskan dengan $\binom{n+k-1}{k-1}$. Untuk $n=15$ dan $k=4$, diperoleh $\binom{15+4-1}{4-1} = \binom{18}{3} = \frac{18 \times 17 \times 16}{3 \times 2 \times 1} = 816$.

\item Banyaknya angka nol berurutan di akhir faktorial ekuivalen dengan mencari pangkat faktor 5 dalam dekomposisi prima nilai tersebut. Menggunakan Teorema Legendre, pangkat dari 5 pada $100!$ adalah $\lfloor \frac{100}{5} \rfloor + \lfloor \frac{100}{25} \rfloor = 20 + 4 = 24$.

\item Misalkan $y = x^2$, persamaan menjadi $y^2 - 5y + 4 = 0$. Bentuk ini dapat difaktorkan menjadi $(y-1)(y-4) = 0$, sehingga solusinya $y = 1$ atau $y = 4$. Substitusi kembali menghasilkan $x^2 = 1 \implies x = \pm 1$ dan $x^2 = 4 \implies x = \pm 2$. Ada 4 akar real yang memenuhi.

\item Total segitiga yang dapat dibentuk dari dodekagon adalah $\binom{12}{3} = 220$. 
Segitiga dengan tepat 1 sisi bersekutu dibentuk dari 1 sisi segi-12 (12 cara) dan 1 titik yang tidak bersebelahan ($12-4=8$ titik), total $12 \times 8 = 96$. 
Segitiga dengan 2 sisi bersekutu sama dengan jumlah titik sudut segi-12, yaitu 12. 
Maka banyaknya segitiga tanpa sisi bersekutu adalah $220 - 96 - 12 = 112$.

\item Kasus ini adalah masalah \textit{derangement} (kekacauan) untuk 5 objek, yaitu menyusun objek sedemikian rupa sehingga tidak ada satupun yang berada di posisi aslinya. Rumus derangement untuk $n=5$ adalah $D_5 = 5! \left( \frac{1}{2!} - \frac{1}{3!} + \frac{1}{4!} - \frac{1}{5!} \right) = 120 \left( \frac{1}{2} - \frac{1}{6} + \frac{1}{24} - \frac{1}{120} \right) = 60 - 20 + 5 - 1 = 44$.

\item Faktorisasi prima dari $10^5$ adalah $2^5 \times 5^5$. Agar suatu pembagi bukan merupakan kelipatan 10, pembagi tersebut tidak boleh memuat faktor 2 dan 5 secara bersamaan. Maka, faktor tersebut harus berbentuk $2^a$ ($0 \le a \le 5$, ada 6 kemungkinan) atau $5^b$ ($0 \le b \le 5$, ada 6 kemungkinan). Karena angka $1$ ($2^0 5^0$) dihitung dua kali, jumlah total faktornya adalah $6 + 6 - 1 = 11$.

\item Total himpunan bagian dari himpunan beranggota 8 elemen adalah $2^8 = 256$. Himpunan bagian yang tidak memenuhi syarat (tidak memiliki bilangan genap) adalah subset yang hanya disusun dari anggota bilangan ganjil $\{1, 3, 5, 7\}$, yang banyaknya $2^4 = 16$. Maka himpunan bagian yang memiliki setidaknya satu angka genap adalah $256 - 16 = 240$.

\item Dari persamaan pertama, substitusi $y = 2 - x$ ke persamaan kedua sehingga diperoleh $x(2-x) - z^2 = 1 \implies 2x - x^2 - z^2 = 1$. Pindah ruaskan persamaan ini membentuk $(x-1)^2 + z^2 = 0$. Karena jumlahan dari dua bilangan kuadrat real yang menghasilkan nol hanya mungkin jika keduanya nol, maka $x-1 = 0 \implies x=1$ dan $z=0$. Nilai $y = 2-1 = 1$. Hanya ada 1 tripel solusi, yaitu $(1, 1, 0)$.

\item Pada prisma segi-$n$ beraturan, sebuah diagonal ruang dibentuk dengan menghubungkan sebuah titik di alas atas dengan sebuah titik di alas bawah yang bukan merupakan titik tepat di bawahnya, dan bukan pula 2 titik yang berada di sebelahnya. Dari setiap titik di alas atas, terdapat $(n - 3)$ pilihan titik di alas bawah. Untuk prisma segi-8, terdapat $8 \times (8-3) = 40$ diagonal ruang.

\item Banyaknya seluruh rute terpendek di bidang Kartesius dari $(0,0)$ menuju $(5,5)$ dihitung dengan kombinasi $\binom{5+5}{5} = \binom{10}{5} = 252$. Banyaknya rute yang diwajibkan melewati $(2,2)$ adalah kombinasi dari $(0,0) \to (2,2)$ dikali kombinasi $(2,2) \to (5,5)$, yaitu $\binom{4}{2} \times \binom{6}{3} = 6 \times 20 = 120$. Maka rute yang tidak melewati titik $(2,2)$ berjumlah $252 - 120 = 132$.

\end{enumerate}
\end{document}
